% !TeX root = ../../Book3.tex
% 纲要标题：### 14.2 微分模的正合列与基变换
% 纲要位置：工作记录/计划纲要.md，第 2410 行。
% 纲要细目（写作范围）：
% * 商环的余法正合列
% * 方程微分给出的展示；尖点微分模的非零挠元及其零化理想
% * 复合环映射的微分序列
% * 局部化及基变换

\section{微分模的正合列与基变换}
\label{b3:sec:1402}


给定代数的生成元以后，导子的值不能任意选择：它必须保持方程的微分关系。两条基本正合列分别处理“添加关系”和“改变基环”，局部化与基变换公式则使这些计算能够转到所需的局部环或基域。

\subsection{商环的余法正合列}
\label{b3:subsec:1402:01}

商映射 \(B\to B/J\) 把 \(f\in J\) 宣布为零，因而微分中也必须添加 \(\mathrm{d}f=0\) 的关系。两条方程乘积的微分仍含方程因子，转到 \(B/J\) 后就为零，所以 \(J^2\) 中的元素不会提供新的这一阶关系；记录这些方程类的是 \(J/J^2\)。

\begin{theorem}[余法正合列]\label{b3:thm:conormal-sequence}
设 \(A\) 为交换环，\(B\) 为交换 \(A\) 代数、\(J\subseteq B\) 为理想、\(C=B/J\)。有 \(C\) 模正合列
\[
J/J^2\xrightarrow{\delta}C\otimes_B\Omega_{B/A}
\longrightarrow\Omega_{C/A}\longrightarrow0,
\qquad \delta(f+J^2)=1\otimes \mathrm{d}f.
\]
模 \(J/J^2\) 称为相应商映射的\term[yufa-mo]{余法模}[conormal module]。映射 \(\delta\) 一般不单射。
\end{theorem}
\begin{proof}
若 \(f\in J^2\)，将它写成有限和 \(\sum_i u_iv_i\)、\(u_i,v_i\in J\)，则 \(\mathrm{d}f=\sum_i(u_i\,\mathrm{d}v_i+v_i\,\mathrm{d}u_i)\) 在张量到 \(C\) 后为零，故 \(\delta\) 良定义。又 \(\mathrm{d}(bf)=b\,\mathrm{d}f+f\,\mathrm{d}b\)，第二项在 \(C\) 上为零，所以 \(\delta\) 是 \(C\) 线性的。

将 \(C\otimes_B\Omega_{B/A}\) 模掉全部 \(1\otimes \mathrm{d}f\)、\(f\in J\)，记商为 \(Q\)。映射 \(b+J\mapsto[1\otimes \mathrm{d}b]\) 良定义，因为改变代表元只增加一个被模掉的微分。它是 \(A\) 导子。若 \(D:C\rightarrow M\) 是任意 \(A\) 导子，复合 \(B\rightarrow C\) 后的导子对应一个 \(C\otimes_B\Omega_{B/A}\rightarrow M\)，且零化各 \(\mathrm{d}f\)，故唯一通过 \(Q\)。于是 \(Q\) 有 \(\Omega_{C/A}\) 的泛性质，得到所需正合列。
\end{proof}

\begin{corollary}\label{b3:cor:differential-presentation}
设 \(A\) 为交换环，\(f_1,\ldots,f_s\in A[X_1,\ldots,X_n]\)，\(B=A[X_1,\ldots,X_n]/(f_1,\ldots,f_s)\)，\(x_i\) 为变量的像。则有展示
\[
B^s\xrightarrow{J_f^{\mathsf t}} B^n\longrightarrow\Omega_{B/A}\longrightarrow0,
\qquad (J_f)_{ij}=\dfrac{\partial f_i}{\partial X_j}(x).
\]
因此有限展示交换代数的微分模为有限展示模。
特别地，若 \(k\) 为特征不等于 \(2,3\) 的域，取 \(B=k[x,y]/(y^2-x^3)\)，并以 \(x,y\) 兼记变量在商中的像，则
\[
\omega=2x\,\mathrm{d}y-3y\,\mathrm{d}x
\]
是 \(\Omega_{B/k}\) 中的非零挠元，且 \(y\omega=x^2\omega=0\)。
\end{corollary}
\begin{proof}
对 \(P=A[X]\) 的商应用余法正合列，\(\Omega_{P/A}\) 为以 \(\mathrm{d}X_i\) 为基的自由模。若 \(f=\sum_i g_if_i\)，在 \(\Omega_{P/A}\) 张量到 \(B\) 后有 \(\mathrm{d}f=\sum_i g_i\,\mathrm{d}f_i\)，因为含 \(f_i\,\mathrm{d}g_i\) 的项消失。因此全部微分关系由各 \(\mathrm{d}f_i\) 生成。映射 \(B^s\to B^n\) 把第 \(i\) 个标准基向量送到 \(\mathrm{d}f_i\) 的坐标列 \(\bigl(\frac{\partial f_i}{\partial X_1}(x),\ldots,\frac{\partial f_i}{\partial X_n}(x)\bigr)^{\mathsf t}\)。通常 Jacobian 矩阵将第 \(i\) 个方程的偏导放在第 \(i\) 行，这里将它们放在第 \(i\) 列，所以所需矩阵是其转置。

在尖点的情形中，上述展示给出
\[
\Omega_{B/k}=(B\,\mathrm{d}x\oplus B\,\mathrm{d}y)/B(2y\,\mathrm{d}y-3x^2\,\mathrm{d}x).
\]
由 \(y^2=x^3\)，有
\[
y\omega=x(2y\,\mathrm{d}y-3x^2\,\mathrm{d}x)=0,
\qquad x^2\omega=y(2y\,\mathrm{d}y-3x^2\,\mathrm{d}x)=0.
\]
为证 \(\omega\ne0\)，由\cref{b3:prop:cusp-normalization}将 \(B\) 识别为 \(k[t^2,t^3]\subseteq k[t]\)，其中 \(x=t^2\)、\(y=t^3\)。若 \(\omega=0\)，它在自由模中的代表必须是关系生成元的一个 \(B\) 倍数，因此存在 \(h\in B\) 满足
\[
-3y=-3hx^2,\qquad 2x=2hy.
\]
由于 \(2\) 在 \(k\) 中可逆，后一等式在分式域 \(k(t)\) 中迫使 \(h=t^{-1}\)，但 \(B\subseteq k[t]\) 不含此元素，矛盾。因此 \(\omega\) 非零，而 \(y\ne0\) 是整环 \(B\) 的非零元素，所以上面的零化等式也证明它是挠元。
\end{proof}

余法映射不单射，表示非零的方程类也可能没有留下非零的微分关系。例如在特征 \(p\) 的域上取 \(B=k[t]\)、\(J=(t^p)\)。由于 \(t^p\notin(t^{2p})\)，它在 \(J/J^2\) 中的类非零，却满足 \(\delta(t^p)=p t^{p-1}\mathrm{d}t=0\)。这一类生成整个余法模，所以此例的余法映射甚至为零。

\subsection{复合环映射的微分序列}
\label{b3:subsec:1402:02}

对 \(A\to B\to C\)，从相对于 \(A\) 求导改为相对于 \(B\) 求导，意味着将 \(B\) 的像也视为常数。因此要在 \(\Omega_{C/A}\) 中进一步模掉这些元素的微分，末端微分模就成为中间微分模的商。

\begin{theorem}\label{b3:thm:differential-transitivity}
设 \(A\rightarrow B\xrightarrow{g}C\) 为交换环同态，则有 \(C\) 模正合列
\[
C\otimes_B\Omega_{B/A}\xrightarrow{\alpha}\Omega_{C/A}
\xrightarrow{\beta}\Omega_{C/B}\longrightarrow0,
\]
其中 \(\alpha(c\otimes \mathrm{d}b)=c\,\mathrm{d}\bigl(g(b)\bigr)\)、\(\beta(\mathrm{d}_{C/A}c)=\mathrm{d}_{C/B}c\)。
\end{theorem}
\begin{proof}
复合 \(B\rightarrow C\xrightarrow{\mathrm{d}_{C/A}}\Omega_{C/A}\) 是 \(A\) 导子，故给出 \(\alpha\)。在 \(\Omega_{C/A}\) 中模掉各 \(\mathrm{d}\bigl(g(b)\bigr)\) 生成的子模后，普遍导子就零化 \(B\) 的像，因此成为 \(B\) 导子。任意 \(B\) 导子首先是 \(A\) 导子，并零化这些指定元素，所以唯一经过此商。该商即 \(\Omega_{C/B}\)，证明正合性。
\end{proof}

首箭头可能不单射。例如 \(A=k\)、\(B=k[t^p]\subseteq C=k[t]\)、\(\operatorname{char}k=p\) 时，\(t^p\) 是 \(B\) 的一个多项式生成元，所以 \(\mathrm{d}_{B/k}(t^p)\ne0\)；但它映到 \(\Omega_{C/k}\) 中时变成 \(\mathrm{d}_{C/k}(t^p)=0\)。

\subsection{微分模的局部化与基变换}
\label{b3:subsec:1402:03}

只在上环 \(B\) 中倒置乘法集 \(S\) 时，基环仍为 \(A\)。分母 \(s\in S\) 的微分一般不为零，所以分式的微分需要包含 \(\mathrm{d}s\) 项。若分母来自基环中的乘法集 \(T\)，它们本来就必须被零化；同步倒置基环与上环后，这些分母的逆也成为常数，此时相对基环是 \(T^{-1}A\)。

\begin{proposition}\label{b3:prop:differential-localization}
设 \(\varphi:A\to B\) 为交换环同态，\(S\subseteq B\) 为乘法集。则有 \(S^{-1}B\) 模自然同构
\[
\Omega_{S^{-1}B/A}\cong S^{-1}\Omega_{B/A},
\qquad \mathrm{d}(b/s)=\dfrac{\mathrm{d}b}{s}-\dfrac{b\,\mathrm{d}s}{s^2}.
\]
其中 \(b\in B\)、\(s\in S\)。
若 \(T\subseteq A\) 为乘法集，记 \(T^{-1}B\) 为 \(B\) 对 \(\varphi(T)\) 的局部化，则还有 \(T^{-1}B\) 模自然同构
\[
\Omega_{T^{-1}B/T^{-1}A}\cong T^{-1}\Omega_{B/A}.
\]
右侧按 \(A\) 在 \(\Omega_{B/A}\) 上的作用局部化。
\end{proposition}
\begin{proof}
对取值于任意 \(S^{-1}B\) 模 \(M\) 的 \(A\) 导子 \(D:B\rightarrow M\)，若它能延拓为 \(\widetilde D:S^{-1}B\to M\)，则等式 \(\widetilde D(s\cdot s^{-1})=0\) 迫使其满足上述分式公式。按此公式定义延拓，检查良定义：若 \(b/s=b'/s'\)，则存在 \(t\in S\) 使 \(t(s'b-sb')=0\)。取导数并在 \(M\) 中将 \(t,s,s'\) 视为单位；由于 \(s'b-sb'=0\) 在局部化后成立，得到
\[
s'Db+bDs'-sDb'-b'Ds=0.
\]
连同 \(s'b=sb'\)，整理即得两个分式公式相等。对分式加法和乘法展开可得加法性与 Leibniz 法则：乘法中的导数项恰分成对第一个和第二个因子的两部分。故延拓存在且唯一，由微分模的泛性质得到第一个同构。

对于同步局部化，取任意 \(T^{-1}B\) 模 \(M\)。每个 \(A\) 导子 \(D:B\to M\) 同样唯一延拓为 \(\widetilde D:T^{-1}B\to M\)。对 \(t\in T\)，有
\[
\widetilde D\bigl(\varphi(t)^{-1}\bigr)
=-\varphi(t)^{-2}\cdot D\bigl(\varphi(t)\bigr)=0.
\]
它又零化 \(\varphi(A)\)，所以零化 \(T^{-1}A\) 的像，因而是 \(T^{-1}A\) 导子。反过来，限制一个 \(T^{-1}A\) 导子便得到 \(A\) 导子。两种操作互逆，再用微分模的泛性质得到第二个同构。
\end{proof}

\begin{proposition}\label{b3:prop:differential-base-change}
设 \(B\) 为交换 \(A\) 代数，\(A\to A'\) 为交换环同态，令 \(B'=A'\otimes_AB\)。有自然同构
\[
\Omega_{B'/A'}\cong B'\otimes_B\Omega_{B/A}.
\]
不需要 \(A'\) 在 \(A\) 上平坦。
\end{proposition}
\begin{proof}
任取 \(B'\) 模 \(M\)，沿环同态 \(B\to B'\)、\(b\mapsto1\otimes b\) 将它看作 \(B\) 模。将 \(A'\) 导子 \(\Delta:B'\to M\) 与这个环同态复合，便给出 \(A\) 导子 \(D:B\to M\)。反过来，给定这样的 \(D\)，在纯张量上定义
\[
\Delta(a'\otimes b)=(a'\otimes1)\cdot D(b).
\]
由 \(D\) 的 \(A\) 线性及张量积中 \(A\) 的两个像相等，此公式保持平衡关系，因而对有限和良定义。它零化 \(A'\) 的像；对纯张量的乘积使用 \(D\) 的 Leibniz 法则，再由加法展开，就得到 \(\Delta\) 的 Leibniz 法则。复合与延拓互逆，因为任何 \(A'\) 导子都满足所列纯张量公式。因此两种导子自然一一对应。

又由普遍导子的泛性质，\(D\) 唯一写成 \(h\circ\mathrm{d}_{B/A}\)，其中 \(h:\Omega_{B/A}\to M\) 为 \(B\) 线性映射。这样的 \(h\) 又唯一延拓为 \(B'\) 线性映射
\[
B'\otimes_B\Omega_{B/A}\longrightarrow M,
\qquad b'\otimes\eta\longmapsto b'\cdot h(\eta).
\]
所以右侧张量积和 \(\Omega_{B'/A'}\) 对每个目标模都表示同一个导子函子，普遍对象的唯一性给出所需同构。这个同构将 \(1\otimes\mathrm{d}b\) 送到 \(\mathrm{d}(1\otimes b)\)。
\end{proof}

这两个公式允许先在多项式商上写出 Jacobian 展示，再局部化到指定点，或改变基域后重新考察方程。代数作为环的约化性和正则性可能在基变换后改变；光滑性则是相对于基环的性质，并在任意基变换下保持，具体论证见\secref{b3:sec:1405}。
